% GENERATED FILE. Edit canonical lessons, then run npm run generate:books.
% canonical-source-hash: fnv1a64:435c9a5a9bf2345c
% canonical-lessons: ES-C433-portal, ES-C433-portera, ES-C433-garaje, ES-C433-norma, ES-C433-junto, ES-R433-repaso-edificio, ES-C433-sintesis-llave

\chapter{El Edificio Y Sus Normas}
\label{ch:es-edificio}
\hlchaptermodality{433}

\begin{chapteropening}
\textbf{By the end of this chapter:} \emph{I can read the notices in a block of flats, say where to leave a key, and say where something is in relation to something else.}

Read the notice or advert this chapter's words exist for, then say the same thing back.
\end{chapteropening}

\section[el portal]{el portal — the big door, and the space behind it}

\label{lesson:ES-C433-portal}



\begin{quote}
Say \emph{la puerta}. Say \emph{el edificio}. Between the street and the stairs of a Spanish block there is a space you have no word for. Here it is.
\end{quote}

\subsection*{You'll want to know: el portal}
\textbf{el portal} — the street door of a block of flats, and the hall behind it.

One word covers both the door and the space. \emph{Te espero en el portal} means downstairs, inside, out of the rain.

\begin{cousinweb}
\textbf{puerta} plus \textbf{-al}, an ending that makes a bigger or more general thing of the noun it lands on.

\noindent
\begin{tabularx}{\linewidth}{@{}>{\raggedright\arraybackslash}X>{\raggedright\arraybackslash}X@{}}
\toprule
\textbf{the thing} & \textbf{the larger thing} \\
\midrule
la puerta & el portal \\
el árbol & la arboleda \\
la rama & el ramal \\
\bottomrule
\end{tabularx}

English has no single word here, which is why translations reach for \emph{doorway}, \emph{entrance}, \emph{lobby} or \emph{hallway} and none quite fits. When a Spanish word has no clean English partner, learning the \textbf{situation} beats learning a gloss: the portal is where you buzz, where the post boxes are, and where you wait.
\end{cousinweb}

\subsection*{Your turn}
Say these aloud:

\begin{itemize}
\raggedright
  \item \emph{el portal}
  \item the door, then the entrance hall
  \item I will wait for you in the portal
\end{itemize}

\begin{tcolorbox}[breakable,colback=teal!4,colframe=teal!35!black,title={Before you move on}]
The street door of a block? (\textbf{el portal}.) Which word is inside it? (\textbf{puerta}.) Does it name the door or the hall? (\textbf{Both.})
\end{tcolorbox}
\section[la portera]{la portera — the person of the portal}

\label{lesson:ES-C433-portera}



\begin{quote}
Say \emph{el portal}. Say \emph{el vecino}. There is somebody who is neither a neighbour nor a stranger, and who is in the portal most mornings.
\end{quote}

\subsection*{You'll want to know: la portera}
\textbf{la portera} — the concierge, the caretaker. \textbf{el portero} if a man.

She cleans the stairs, takes in parcels, knows who is away, and is the person a notice tells you to leave a key with.

\begin{cousinweb}
\textbf{portal} and \textbf{portera} share \textbf{puerta}, and the ending \textbf{-ero / -era} makes the \textbf{person associated with} a thing:

\noindent
\begin{tabularx}{\linewidth}{@{}>{\raggedright\arraybackslash}X>{\raggedright\arraybackslash}X@{}}
\toprule
\textbf{the thing} & \textbf{the person} \\
\midrule
la puerta & la portera \\
el pan & el panadero \\
el zapato & el zapatero \\
\bottomrule
\end{tabularx}

You met \textbf{-ería} for the place and \textbf{-dor} for the doer. This is the third: \textbf{-ero} for the person whose job is that thing. Three endings, three different questions, and all of them turn nouns you have into words you have not met.

\textbf{el portero} also means the goalkeeper. The one who guards the gate, in both cases.
\end{cousinweb}

\subsection*{Your turn}
Say these aloud:

\begin{itemize}
\raggedright
  \item \emph{la portera}
  \item leave the key with the concierge
  \item the person who works with bread, using the ending
\end{itemize}

\begin{tcolorbox}[breakable,colback=teal!4,colframe=teal!35!black,title={Before you move on}]
The concierge? (\textbf{la portera}.) What does \textbf{-ero} name? (\textbf{The person whose job it is.}) What else is \emph{el portero}? (\textbf{The goalkeeper.})
\end{tcolorbox}
\section[el garaje]{el garaje — spelled to be said}

\label{lesson:ES-C433-garaje}



\begin{quote}
Say \emph{el coche}. Say \emph{la portera}. Now name the floor below the portal where the cars live.
\end{quote}

\subsection*{You'll want to know: el garaje}
\textbf{el garaje} — the garage, the underground car park of a building.

Said \emph{ga-RA-je}, with the \textbf{j} as the throaty \emph{jota}.

\begin{cousinweb}
French \textbf{garage} came into Spanish and was \textbf{respelled}. French writes the sound at the end with \emph{-ge}; Spanish has no such spelling, so it wrote what it heard with the letter it uses for that sound: \textbf{j}.

\noindent
\begin{tabularx}{\linewidth}{@{}>{\raggedright\arraybackslash}X>{\raggedright\arraybackslash}X@{}}
\toprule
\textbf{language} & \textbf{spelling} \\
\midrule
French & garage \\
Spanish & garaje \\
\bottomrule
\end{tabularx}

That habit is worth more than the word. Spanish borrows freely but \textbf{respells to its own rules}, which is why \emph{fútbol}, \emph{champú} and \emph{chófer} look the way they do. A borrowed word you meet in Spanish has usually been made to spell itself the Spanish way, so read it aloud by Spanish rules and the original often appears.
\end{cousinweb}

\subsection*{Your turn}
Say these aloud:

\begin{itemize}
\raggedright
  \item \emph{el garaje}
  \item the car is in the garage
  \item why it is not spelled the French way
\end{itemize}

\begin{tcolorbox}[breakable,colback=teal!4,colframe=teal!35!black,title={Before you move on}]
The garage? (\textbf{el garaje}.) Which language did it come from? (\textbf{French}.) What did Spanish change? (\textbf{The spelling.})
\end{tcolorbox}
\section[la norma]{la norma — what is written on the wall}

\label{lesson:ES-C433-norma}



\begin{quote}
Say \emph{el portal}. Say \emph{el garaje}. Both have a list on the wall. Name what that list is called.
\end{quote}

\subsection*{You'll want to know: la norma}
\textbf{la norma} — the rule. \textbf{las normas} — the rules, the regulations.

You meet it in the plural almost always, as a heading: \emph{NORMAS DE USO}, \emph{NORMAS DE PRÉSTAMO}.

\begin{cousinweb}
Latin \textbf{norma} was a carpenter's square — the tool for checking a right angle. A thing that matched it was \emph{correct}; a thing that did not was off. The word travelled from the tool to the standard itself, and from there into every European language.

That history explains a family you can now read at sight:

\noindent
\begin{tabularx}{\linewidth}{@{}>{\raggedright\arraybackslash}X>{\raggedright\arraybackslash}X@{}}
\toprule
\textbf{word} & \textbf{what it holds} \\
\midrule
la norma & the standard \\
normal & matching the standard \\
enorme & out of the square, huge \\
\bottomrule
\end{tabularx}

\textbf{enorme} is the one worth noticing. \emph{E-} out of, plus \emph{norma}: outside the square. Something enormous is something that will not fit the measure.
\end{cousinweb}

\subsection*{Your turn}
Say these aloud:

\begin{itemize}
\raggedright
  \item \emph{las normas}
  \item the rules of the building
  \item what \emph{enorme} has to do with \emph{norma}
\end{itemize}

\begin{tcolorbox}[breakable,colback=teal!4,colframe=teal!35!black,title={Before you move on}]
The rules? (\textbf{las normas}.) What was a Latin \emph{norma}? (\textbf{A carpenter's square}.) And \emph{enorme}? (\textbf{Outside the square.})
\end{tcolorbox}
\section[junto]{junto — together, or beside}

\label{lesson:ES-C433-junto}



\begin{quote}
Say \emph{la calle}. You can put a thing \textbf{in} a street. Say that two people did something \textbf{together} and see whether you have the word.
\end{quote}

\subsection*{You'll want to know: junto}
\textbf{juntos} — together. \textbf{junto a} — next to, beside.

Same word, two jobs, and the little \textbf{a} is what separates them.

\begin{grammarlens}[title={Grammar Lens: the a decides}]
\noindent
\begin{tabularx}{\linewidth}{@{}>{\raggedright\arraybackslash}X>{\raggedright\arraybackslash}X@{}}
\toprule
\textbf{you say} & \textbf{it means} \\
\midrule
cenamos juntos & we had dinner together \\
junto a la puerta & next to the door \\
\bottomrule
\end{tabularx}

Without \textbf{a}, it describes the people and agrees with them: \emph{juntos}, \emph{juntas}. With \textbf{a}, it stops describing anybody and becomes a position — and then it stops agreeing, staying \textbf{junto} whatever follows.

That is a small thing to watch, and it is the sort of small thing a reader trips on: \emph{las cajas, juntas} is a comment about the boxes; \emph{junto a las cajas} is where you are standing.
\end{grammarlens}

\subsection*{Your turn}
Say these aloud:

\begin{itemize}
\raggedright
  \item we went together
  \item next to the portal
  \item which of the two agrees with the people
\end{itemize}

\begin{tcolorbox}[breakable,colback=teal!4,colframe=teal!35!black,title={Before you move on}]
Together? (\textbf{juntos}.) Next to? (\textbf{junto a}.) Which one changes for gender? (\textbf{The one without \emph{a}.})
\end{tcolorbox}
\section[(portal, portera, garaje, norma, junto)]{(portal, portera, garaje, norma, junto) — from cold}

\label{lesson:ES-R433-repaso-edificio}



\begin{quote}
Close the lessons before this one. Say \emph{the entrance hall}, \emph{the concierge} and \emph{the garage}.
\end{quote}

\begin{grammarlens}[title={Grammar Lens: three endings, three questions}]
\noindent
\begin{tabularx}{\linewidth}{@{}>{\raggedright\arraybackslash}X>{\raggedright\arraybackslash}X@{}}
\toprule
\textbf{ending} & \textbf{the question it answers} \\
\midrule
\textbf{-dor} & who or what does it? \\
\textbf{-ería} & where is it kept? \\
\textbf{-ero} & whose job is it? \\
\bottomrule
\end{tabularx}

All three land on \textbf{puerta} or a word like it and hand you a noun you never learned. Do not read a rule into it — many words have no ending at all, and \emph{el garaje} and \emph{la norma} in this very chapter came in whole from elsewhere.

What these three chapters have built is a \textbf{habit}: before treating a long word as new, look for a piece you already hold.
\end{grammarlens}

\begin{grammarlens}[title={What you've built}]
\textbf{You could describe a home. You could not describe a building.} \emph{La casa}, \emph{la puerta} and \emph{el vecino} were yours; the hall, the person who cleans it, the car park beneath it and the list of rules on its wall were not.

Those four are where written Spanish actually meets a resident — on a notice, in a message from an administrator, in an instruction about a key. And \textbf{junto a} is the small word that lets you say where any of it is.
\end{grammarlens}

\subsection*{Your turn}
Say these aloud:

\begin{itemize}
\raggedright
  \item all five words
  \item leave the key with the concierge, next to the door
  \item the rules of the garage
  \item the three endings and what each asks
\end{itemize}

\begin{tcolorbox}[breakable,colback=teal!4,colframe=teal!35!black,title={Before you move on}]
The concierge? (\textbf{la portera}.) The rules? (\textbf{las normas}.) Next to? (\textbf{junto a}.)
\end{tcolorbox}
\section[Practice]{(dejar una llave) — read it, then do what it says}

\label{lesson:ES-C433-sintesis-llave}



\begin{quote}
Say \emph{la portera} and \emph{el contador}. Both appear in the notice you are about to read, and the notice asks you to do something.
\end{quote}

\subsection*{Your turn}
Read it once straight through, then answer.

\begin{quote}
\textbf{NORMAS DEL PORTAL} El jueves vendrán a leer los \textbf{contadores} entre las nueve y las dos. Es necesario que haya alguien en casa. Quien no pueda estar debe dejar una llave a la \textbf{portera}.
\end{quote}

Say these aloud:

\begin{itemize}
\raggedright
  \item what day they are coming and what for
  \item what you must do if you cannot be at home
  \item who you leave the key with
\end{itemize}

\begin{grammarlens}[title={What you've built}]
Now produce it yourself. Say aloud:

Say these aloud:

\begin{itemize}
\raggedright
  \item I cannot be at home on Thursday
  \item I am leaving the key with the concierge
  \item my car is in the garage, next to the door
\end{itemize}

\textbf{This is the level's real test, and it is not a conversation.} A2 asks you to read a short official notice and act on it correctly. Three chapters ago the words for the notice did not exist for you; the grammar never changed.
\end{grammarlens}

\begin{tcolorbox}[breakable,colback=teal!4,colframe=teal!35!black,title={Before you move on}]
Who do you leave the key with? (\textbf{la portera}.) Where do you wait? (\textbf{en el portal}.) What is the heading on the wall? (\textbf{NORMAS}.)
\end{tcolorbox}
